% SPDX-License-Identifier: MIT
% Copyright (c) 2026 Evgeny Baulin
% Optimization Methods, Seminar 08: theory handout (English).
% Generated by checks/08. Constrained Optimality and Discrete Modelling.ipynb from web/08/js/content-*.js and content/students.json. Do not edit by hand.
% Compile with: latexmk -xelatex -interaction=nonstopmode -halt-on-error Theory_en.tex
\documentclass[11pt,a4paper]{article}

\usepackage{fontspec}
\setmainfont{cmunrm.otf}[BoldFont=cmunbx.otf, ItalicFont=cmunti.otf, BoldItalicFont=cmunbi.otf]
\setsansfont{cmunss.otf}[BoldFont=cmunsx.otf, ItalicFont=cmunsi.otf, BoldItalicFont=cmunso.otf]
\setmonofont{cmuntt.otf}[BoldFont=cmuntb.otf, ItalicFont=cmunit.otf, BoldItalicFont=cmuntx.otf]
\usepackage{polyglossia}
\setdefaultlanguage{english}

\usepackage[a4paper,left=2.3cm,right=2.3cm,top=2.4cm,bottom=2.4cm,headheight=14pt]{geometry}
\usepackage{amsmath,amssymb}
\usepackage{graphicx}
\usepackage{array,booktabs,longtable}
\usepackage{enumitem}
\usepackage{xcolor}
\usepackage[most]{tcolorbox}
\usepackage{titlesec}
\usepackage{fancyhdr}
\usepackage{needspace}
\usepackage[expansion=false]{microtype}
\usepackage[colorlinks=true,linkcolor=hseblue,urlcolor=hseblue,citecolor=hseblue]{hyperref}

\definecolor{hseblue}{RGB}{16,45,105}
\definecolor{hseaccent}{RGB}{198,40,40}
\definecolor{hsegray}{RGB}{108,114,126}
\definecolor{hselight}{RGB}{240,243,248}
\definecolor{hsewarm}{RGB}{253,244,232}

\hypersetup{pdftitle={Optimization Methods. Seminar 08. Theory handout}, pdfauthor={Evgeny Baulin}}

\titleformat{\section}{\color{hseblue}\Large\bfseries}{\thesection}{0.6em}{}
\titleformat{\subsection}{\color{hseblue}\large\bfseries}{\thesubsection}{0.6em}{}
\titlespacing*{\section}{0pt}{16pt plus 4pt minus 2pt}{8pt plus 2pt}
\titlespacing*{\subsection}{0pt}{12pt plus 3pt minus 2pt}{6pt plus 2pt}
\pretocmd{\section}{\Needspace{7\baselineskip}}{}{\errmessage{cannot patch section}}
\pretocmd{\subsection}{\Needspace{5\baselineskip}}{}{\errmessage{cannot patch subsection}}

\tcbset{handoutbox/.style={enhanced, breakable, sharp corners, boxrule=0pt, left=7pt, right=7pt, top=5pt, bottom=5pt, before skip=8pt plus 2pt, after skip=8pt plus 2pt}}
\newtcolorbox{examplebox}{handoutbox, colback=hselight, colframe=hselight, before upper={\textcolor{hseblue}{\bfseries Example.}\ }}
\newtcolorbox{whybox}{handoutbox, colback=white, colframe=hseblue, leftrule=2.5pt, top=3pt, bottom=3pt, before upper={\textcolor{hseaccent}{\bfseries Why.}\ }}
\newtcolorbox{answerbox}{handoutbox, colback=hsewarm, colframe=hsewarm, before upper={\textcolor{hseaccent}{\bfseries Answer.}\ }}
\newtcolorbox{conventionbox}{handoutbox, colback=hselight, colframe=hseblue, leftrule=2.5pt, before upper={\textcolor{hseblue}{\bfseries Notation.}\par\smallskip}}

\newcommand{\handoutfigure}[3]{\par\medskip\noindent\begin{minipage}{\linewidth}\centering
  \includegraphics[width=#1\linewidth]{figures/#2.pdf}\par\smallskip
  \begin{minipage}{0.94\linewidth}\small\textcolor{hseblue}{\bfseries Figure.}\ #3\end{minipage}\end{minipage}\par\medskip}
\newcommand{\problemhead}[3]{\par\bigskip\noindent{\color{hseblue}\large\bfseries #1.\ #2}\par\nopagebreak
  \noindent{\small\itshape\color{hsegray}#3}\par\nopagebreak\smallskip}
\newtcolorbox{definitionbox}[1]{handoutbox, colback=white, colframe=hseblue, leftrule=2.5pt, before upper={\textcolor{hseblue}{\bfseries #1}\ }}
\newtcolorbox{theorembox}[1]{handoutbox, colback=hselight, colframe=hseblue, leftrule=2.5pt, before upper={\textcolor{hseblue}{\bfseries #1}\ }}
\newtcolorbox{proofbox}{handoutbox, colback=white, colframe=hsegray, leftrule=1.2pt, top=3pt, bottom=3pt, before upper={\textcolor{hsegray}{\bfseries\itshape Proof.}\ }}
\newtcolorbox{notebox}[1]{handoutbox, colback=hselight, colframe=hselight, before upper={\textcolor{hseblue}{\bfseries #1}\par\smallskip}}
\newtcolorbox{trapbox}{handoutbox, colback=hsewarm, colframe=hseaccent, leftrule=2.5pt, before upper={\textcolor{hseaccent}{\bfseries Trap.}\ }}
\newtcolorbox{selfcheckbox}{handoutbox, colback=white, colframe=hseaccent, leftrule=2.5pt, top=3pt, bottom=3pt, before upper={\textcolor{hseaccent}{\bfseries Check yourself.}\ }}
\newcommand{\handoutdisplay}[1]{\sbox0{$\displaystyle #1$}\ifdim\wd0>\linewidth
  \par{\centering\binoppenalty=10000\relpenalty=0\relax$\displaystyle #1$\par}\else\[#1\]\fi}
\newcommand{\selfcheckanswer}{\par\smallskip\textcolor{hsegray}{\bfseries Answer.}\ }

\setlist[itemize]{itemsep=3pt,topsep=4pt}
\setlist[enumerate]{itemsep=3pt,topsep=4pt}
\setlength{\parindent}{0pt}
\setlength{\parskip}{5pt plus 1pt}
\setlength{\emergencystretch}{2em}
\setcounter{tocdepth}{2}

\pagestyle{fancy}
\fancyhf{}
\fancyhead[L]{\small\color{hsegray}Optimization Methods}
\fancyhead[R]{\small\color{hsegray}Seminar 08}
\fancyfoot[L]{\small\color{hsegray}Evgeny Baulin}
\fancyfoot[C]{\small\thepage}
\fancyfoot[R]{\small\color{hsegray}MIT License}
\renewcommand{\headrulewidth}{0pt}
\renewcommand{\footrulewidth}{0pt}
\fancypagestyle{plain}{\pagestyle{fancy}}

\begin{document}

\begin{center}
  {\large\color{hsegray}Optimization Methods\par}
  \vspace{4pt}
  {\large\bfseries\color{hseaccent}Seminar 08\par}
  \vspace{8pt}
  {\LARGE\bfseries\color{hseblue}Constrained Optimality and Discrete Modelling: KKT, Duality, Integer Models\par}
  \vspace{8pt}
  {\Large Theory handout\par}
  \vspace{12pt}
  HSE University, Faculty of Computer Science\par
  Academic year 2026-2027\par
  \vspace{8pt}
  Course lecturer: Andrey Ignatov\par
  Evgeny Baulin\par
  Telegram: \href{https://t.me/tarakan_tuc}{@tarakan\_tuc}\par
\end{center}

\noindent{\color{hseblue}\rule{\linewidth}{0.6pt}}

\tableofcontents

\clearpage


This handout follows seminar 8. It goes from the first-order conditions of a constrained problem to the dual problem and the meaning of multipliers, and then turns to models with integer variables, where the same idea of a bound reappears as the relaxation. All nine problems of the class are solved in full.

\begin{conventionbox}
Problem: minimize $f(x)$ over $x\in\mathbb R^n$ subject to $g_i(x)\le0$ and $h_j(x)=0$. Multipliers $\lambda_i\ge0$ for inequalities, $\nu_j$ of any sign for equalities. $L=f+\sum_i\lambda_ig_i+\sum_j\nu_jh_j$, $q=\inf_xL$, $p^*$ and $d^*$ the primal and dual optimal values. In the integer part a maximization keeps its own sign: there the relaxation bounds the optimum from above.
\end{conventionbox}

% =====================================================================
\section{Optimality conditions with constraints}\label{sec:kkt}

An unconstrained minimizer has $\nabla f=0$. With constraints this fails as soon as the solution sits on the boundary: $f$ would still like to go further, but the wall stops it. What remains true is a balance between the pull of $f$ and the push of the walls that are touched.

\begin{definitionbox}{Definition: Active set.}
For a feasible $x$, the indices $i$ with $g_i(x)=0$ form the active set $A(x)$. Constraints with $g_i(x)<0$ are slack and do not restrict small moves.
\end{definitionbox}

\begin{theorembox}{Theorem: Karush–Kuhn–Tucker conditions.}
Let $x^*$ be a local minimizer and suppose a constraint qualification holds there (for instance, the gradients of the active constraints and of the $h_j$ are linearly independent, or the problem is convex and Slater's condition holds). Then there are $\lambda^*$ and $\nu^*$ such that $\nabla_xL(x^*,\lambda^*,\nu^*)=0$; $x^*$ is feasible; $\lambda^*\ge0$; and $\lambda_i^*g_i(x^*)=0$ for every $i$.
\end{theorembox}

\begin{whybox}
Stationarity says $-\nabla f(x^*)=\sum_i\lambda_i^*\nabla g_i(x^*)+\sum_j\nu_j^*\nabla h_j(x^*)$: the direction in which $f$ falls fastest points out of the feasible set, into the cone spanned by the outward normals of the touched walls. Any feasible move then has to go uphill.
\end{whybox}

\begin{theorembox}{Proposition: Sufficiency in the convex case.}
If $f$ and every $g_i$ are convex and every $h_j$ is affine, a feasible point that satisfies the KKT conditions with some multipliers is a global minimizer.
\end{theorembox}

\begin{proofbox}
With $\lambda^*\ge0$ fixed, $x\mapsto L(x,\lambda^*,\nu^*)$ is convex and has zero gradient at $x^*$, so $x^*$ minimizes it over all of $\mathbb R^n$. Complementary slackness and $h(x^*)=0$ give $L(x^*,\lambda^*,\nu^*)=f(x^*)$. For a feasible $x$ the added terms are $\le0$, so $f(x^*)\le L(x,\lambda^*,\nu^*)\le f(x)$.
\end{proofbox}

\begin{itemize}
  \item Write $L$ and the stationarity equations.
  \item Guess an active set $A$. Put $g_i=0$ for $i\in A$ and $\lambda_i=0$ for $i\notin A$, then solve the square system.
  \item Discard the candidate if some $g_i>0$ or some $\lambda_i<0$.
  \item In a convex problem the first survivor is the answer; otherwise compare the values of all survivors and check that a minimizer exists.
\end{itemize}

\handoutfigure{0.62}{fig_a1}{In A1 the minimizer lies on the line of $g_1$; $-\nabla f$ is a positive multiple of the normal of that line.}

\problemhead{A1}{Four guesses for the active set}{On the board}
\textbf{Statement.} Minimize $f(x,y)=(x-3)^2+2(y-2)^2$ subject to $g_1=x+y-2\le0$ and $g_2=-y\le0$. Go through all four active sets, find the KKT point and its multipliers, and give the first coordinate of $x^*$, $\lambda_1^*$ and $f(x^*)$.

\textbf{Solution.}
\begingroup\sloppy
\begin{enumerate}[label=\textbf{\arabic*.}, leftmargin=*]
  \item No active constraint: the unconstrained minimizer $(3,\ 2)$ has $g_1=3>0$, so it is infeasible and at least one constraint binds.
  \item Only $g_2$ active: $y=0$, $\lambda_1=0$, and the first equation gives $x=3$; the point $(3,\ 0)$ has $g_1=1>0$ and is rejected. Both active: $x+y=2$ and $y=0$ give $(2,\ 0)$, and stationarity gives $\lambda=(2,\ -6)$. A negative multiplier means that the floor $y\ge0$ pushes the point the wrong way: rejected.
  \item Only $g_1$ active, $\lambda_2=0$: $x=3-\lambda_1/2$, $y=2-\lambda_1/4$, and $x+y=2$ gives $5-3\lambda_1/4=2$, so $\lambda_1=4$. The candidate is $x^*=(1,\ 1)$ with $\lambda^*=(4,\ 0)$; $y^*>0$, so $g_2$ is indeed inactive.
  \item At $x^*$ the gradient is $\nabla f=(-4,\ -4)=-\lambda_1^*\nabla g_1$: minus the gradient points straight out through the line $x+y=2$. The problem is convex, so the KKT point is the global minimizer, and $f(x^*)=6$.
\end{enumerate}
\endgroup
\begin{answerbox}
$x^*=(1,\ 1)$, $\lambda^*=(4,\ 0)$, $f(x^*)=6$.
\end{answerbox}

\begin{trapbox}
Active does not mean positive multiplier. The implication goes one way: $\lambda_i>0$ forces $g_i=0$, but $g_i=0$ allows $\lambda_i=0$. A constraint that just touches the unconstrained solution is the typical case.
\end{trapbox}

\problemhead{A2}{Three constraints}{On your own}
\textbf{Statement.} Minimize $f(x,y)=3(x-5)^2+(y-2)^2$ subject to $g_1=x-3.5\le0$, $g_2=x+y-5.5\le0$ and $g_3=-y\le0$. Find the KKT point, the active set and all three multipliers. Give $\lambda_1^*$, $\lambda_2^*$ and $f(x^*)$.

\textbf{Solution.}
\begingroup\sloppy
\begin{enumerate}[label=\textbf{\arabic*.}, leftmargin=*]
  \item Without constraints $f$ would sit at $(5,2)$, but there $g_1>0$: the bound $x\le3.5$ is violated, so $g_1$ must be tight. Put $x=3.5$ and keep the other multipliers at zero for a moment. The $y$-equation $2(y-2)=0$ gives $y=2$, the $x$-equation $6(x-5)+\lambda_1=0$ gives $\lambda_1=9$.
  \item At $x^*=(3.5,\ 2)$ the three constraint values are $g=(0,\ 0,\ -2)$. The sum constraint $g_2$ is tight as well, although nobody asked for it, while $g_3$ is slack, so $\lambda_3=0$ by complementary slackness. The active set is $\{1,2\}$.
  \item With both $g_1$ and $g_2$ in play, stationarity reads $6(x-5)+\lambda_1+\lambda_2=0$ and $2(y-2)+\lambda_2-\lambda_3=0$. At $y=2$, $\lambda_3=0$ the second equation forces $\lambda_2=0$, and then the first gives $\lambda_1=9$ again. The gradients $(1,0)$ and $(1,1)$ of the tight constraints are independent, so these multipliers are the only ones: $\lambda^*=(9,\ 0,\ 0)$.
  \item All four conditions hold, $f$ is convex and the constraints are linear, so the point is the global minimizer with $f^*=6.75$. The constraint $g_2$ touches $x^*$ but carries no weight: delete it and nothing moves; delete $g_1$ and the answer changes.
\end{enumerate}
\endgroup
\begin{answerbox}
$x^*=(3.5,\ 2)$, $\lambda_1=9$, $\lambda_2=0$, $\lambda_3=0$, $f^*=6.75$.
\end{answerbox}

\problemhead{A3}{A linear objective on a disk}{On your own}
\textbf{Statement.} Minimize $x+2y$ subject to $x^2+y^2\le45$. Explain why the minimizer lies on the boundary, find it from the KKT conditions and check that it is global. Give $y^*$, $\lambda^*$ and $f(x^*)$.

\textbf{Solution.}
\begingroup\sloppy
\begin{enumerate}[label=\textbf{\arabic*.}, leftmargin=*]
  \item With $c=(1,\ 2)$ and $g=x^2+y^2-45$, stationarity is $c+2\lambda(x,y)=0$, so $(x,y)=-c/(2\lambda)$: the minimizer lies opposite to $c$, and $\lambda>0$ is needed.
  \item The circle is tight: $\|c\|^2/(4\lambda^2)=45$ with $\|c\|^2=5$, so $\lambda^*=0.1667$, $x^*=(-3,\ -6)$ and $f^*=-15$.
  \item The objective is linear and the disk convex, so this KKT point is the global minimizer; the other point of the circle on the same line is the maximizer, with a negative multiplier.
\end{enumerate}
\endgroup
\begin{answerbox}
$x^*=(-3,\ -6)$, $\lambda^*=0.1667$, $f^*=-15$.
\end{answerbox}

\begin{selfcheckbox}
In A1, why is the case with both constraints active rejected, although its point is feasible?
\selfcheckanswer Its multiplier for $y\ge0$ comes out negative, $\lambda_2=-6$: the wall $y=0$ would have to pull the point instead of pushing it.
\end{selfcheckbox}

% =====================================================================
\section{Lagrangian duality and prices}\label{sec:duality}

\begin{definitionbox}{Definition: Dual function and dual problem.}
$q(\lambda,\nu)=\inf_xL(x,\lambda,\nu)$, possibly $-\infty$. The dual problem is to maximize $q$ over $\lambda\ge0$ and all $\nu$; its value is $d^*$.
\end{definitionbox}

For fixed $x$ the Lagrangian is affine in the multipliers, and an infimum of affine functions is concave. So the dual problem is always a concave maximization, whatever $f$ and the constraints look like.

\begin{theorembox}{Theorem: Weak duality.}
For every feasible $x$, every $\lambda\ge0$ and every $\nu$, $q(\lambda,\nu)\le f(x)$. Hence $d^*\le p^*$.
\end{theorembox}

\begin{proofbox}
$q(\lambda,\nu)\le L(x,\lambda,\nu)=f(x)+\sum_i\lambda_ig_i(x)+\sum_j\nu_jh_j(x)\le f(x)$, since $h(x)=0$ and each $\lambda_ig_i(x)\le0$.
\end{proofbox}

\begin{theorembox}{Theorem: Strong duality under Slater's condition.}
If the problem is convex with affine equalities, $p^*$ is finite and some feasible point satisfies every inequality strictly, then $d^*=p^*$ and the dual maximum is attained.
\end{theorembox}

\begin{whybox}
Read $\lambda_i$ as a fine per unit of violation of $g_i\le0$. The dual asks for the fines at which the unconstrained minimizer of the fined objective is already feasible. Too low a fine and the minimizer cheats; the right fine is exactly the market price of the constraint.
\end{whybox}

\begin{theorembox}{Proposition: Multipliers measure sensitivity.}
Write the constraints as $g_i(x)\le b_i$ and let $p^*(b)$ be the optimal value. Where $p^*$ is differentiable and the optimal multipliers are unique, $\partial p^*/\partial b_i=-\lambda_i^*$.
\end{theorembox}

\begin{trapbox}
Watch the direction of the inequality. In B1 the constraint is a lower bound $x+y\ge b$, written as $b-x-y\le0$; raising $b$ tightens it, and the slope of $p^*$ is $+\lambda^*$, not $-\lambda^*$.
\end{trapbox}

\problemhead{B1}{A dual function by hand}{On the board}
\textbf{Statement.} Minimize $x^2+3y^2$ subject to $x+y\ge4$. Write the constraint as $g\le0$, find the dual function $q(\lambda)$, maximize it and compare $d^*$ with $p^*$. Then treat the right-hand side $b$ as a parameter and find $dp^*/db$ at $b=4$. Give $\lambda^*$, $d^*$ and $dp^*/db$.

\textbf{Solution.}
\begingroup\sloppy
\begin{enumerate}[label=\textbf{\arabic*.}, leftmargin=*]
  \item $g=b-x-y\le0$ and $L=x^2+3y^2+\lambda(b-x-y)$. From $2x-\lambda=0$ and $6y-\lambda=0$ the minimizer over $(x,y)$ is $x=\lambda/2$, $y=\lambda/6$.
  \item Substituting, $q(\lambda)=\lambda^2/4+3(\lambda/6)^2+\lambda(4-2\lambda/3)=4\lambda-\lambda^2/3$, a concave parabola. Its maximum over $\lambda\ge0$ is where $q'(\lambda)=4-2\lambda/3=0$: $\lambda^*=6$ and $d^*=12$.
  \item The primal point comes from the same formulas, $(x^*,y^*)=(3,\ 1)$; it satisfies $x+y=4$ and gives $p^*=12=d^*$. There is no gap, as Slater promises: the point $(4,4)$ is strictly feasible.
  \item For a general $b$ the same steps give $\lambda^*=3b/2$ and $p^*(b)=0.75\,b^2$, so $dp^*/db=6$ at $b=4$, which is $\lambda^*$. The sign is plus because $x+y\ge b$ reads $-x-y\le-b$: raising $b$ tightens the constraint, and each unit costs $\lambda^*$.
\end{enumerate}
\endgroup
\begin{answerbox}
$\lambda^*=6$, $d^*=p^*=12$, $dp^*/db=6$.
\end{answerbox}

\problemhead{B2}{What one more unit of the resource is worth}{On your own}
\textbf{Statement.} Minimize $(x-4)^2+(y-3)^2$ subject to $x+2y\le b$ with $b=3$. Find $q(\lambda)$ and its maximizer, then the primal solution. Predict $p^*$ for $b+0.5$ from the multiplier, then compute it exactly. Give $\lambda^*$, $p^*$ and the exact $p^*$ at the new $b$.

\textbf{Solution.}
\begingroup\sloppy
\begin{enumerate}[label=\textbf{\arabic*.}, leftmargin=*]
  \item With $c=(4,\ 3)$ and $a=(1,\ 2)$, $L=\|z-c\|^2+\lambda(a^\top z-b)$ for $z=(x,y)$. Setting the gradient in $z$ to zero gives $z(\lambda)=c-\tfrac{\lambda}{2}a$, that is $x=4-\lambda/2$, $y=3-\lambda$. The centre itself has $a^\top c$ far above $b=3$, so the constraint will be tight.
  \item Substituting $z(\lambda)$ gives the dual function $q(\lambda)=\lambda(a^\top c-b)-\tfrac{1}{4}\|a\|^2\lambda^2$, a downward parabola. Its top is at $\lambda^*=2.8$ with value $9.8$. The same $\lambda$ put into $z(\lambda)$ gives $x^*=(2.6,\ 0.2)$, which satisfies $a^\top x^*=b$, and $p^*=\|x^*-c\|^2=9.8$: the two values agree, as Slater's condition promises (the origin is strictly feasible).
  \item The constraint is an upper bound $a^\top z\le b$, so $dp^*/db=-\lambda^*$. For $\Delta b=0.5$ the first-order forecast is $p^*-\lambda^*\Delta b=8.4$.
  \item Solving again with the larger right-hand side gives $\lambda=2.6$, $x=(2.7,\ 0.4)$ and $p^*=8.45$. The forecast misses by $0.05$, a second-order effect: here $p^*(b)=(a^\top c-b)^2/\|a\|^2$ is a parabola in $b$, and its tangent lies below it.
\end{enumerate}
\endgroup
\begin{answerbox}
$\lambda^*=2.8$, $p^*=d^*=9.8$; after the change of $b$, $p^*=8.45$ (forecast $8.4$).
\end{answerbox}

\problemhead{B3}{An equality and a free multiplier}{On your own}
\textbf{Statement.} Minimize $x^2+2y^2$ subject to $h=x+y-6=0$. Find the dual function $q(\nu)$, maximize it over all $\nu\in\mathbb R$ and recover the primal point. Give $\nu^*$, $d^*$ and $x^*$.

\textbf{Solution.}
\begingroup\sloppy
\begin{enumerate}[label=\textbf{\arabic*.}, leftmargin=*]
  \item $L=x^2+2y^2+\nu(x+y-6)$. Minimizing over $(x,y)$ gives $x=-0.5\nu$, $y=-0.25\nu$.
  \item Substituting, $q(\nu)=-6\nu-0.375\nu^2$, a downward parabola. Its top is at $\nu^*=-8$ with $d^*=24$.
  \item Then $x^*=(4,\ 2)$ lies on the line and $p^*=24=d^*$. The multiplier came out negative, which is allowed for an equality: it only says on which side of the line $f$ would like to be.
\end{enumerate}
\endgroup
\begin{answerbox}
$\nu^*=-8$, $d^*=p^*=24$, $x^*=(4,\ 2)$.
\end{answerbox}

\begin{selfcheckbox}
Can $q(\lambda)$ ever lie above $p^*$?
\selfcheckanswer No, for any $\lambda\ge0$: that is weak duality, and it needs neither convexity nor Slater.
\end{selfcheckbox}

% =====================================================================
\section{From a word problem to an integer model}\label{sec:modelling}

Turning a story into a model takes three passes. Name the decisions and give each a variable, stating whether it is continuous, integer or binary. Translate every rule of the story into a constraint. Finally write the goal as a function of the variables. Most errors come from a rule that was never written down.

\begin{definitionbox}{Definition: Integer program and its relaxation.}
A mixed-integer linear program has a linear objective, linear constraints and some variables restricted to integers. Its LP relaxation is the same problem with the integrality requirement dropped.
\end{definitionbox}

\begin{theorembox}{Proposition: Dropping integrality gives a bound.}
For a maximization the relaxed optimum is at least the integer optimum; for a minimization it is at most. If the relaxed optimum happens to be integer, it solves the integer problem.
\end{theorembox}

\begin{proofbox}
Every integer feasible point is feasible for the relaxation, so the relaxation optimizes over a larger set.
\end{proofbox}

\handoutfigure{0.62}{fig_c1}{The relaxation of C1 peaks at a fractional vertex; the best integer point is not one of its roundings.}

\problemhead{C1}{Chairs and tables}{On the board}
\textbf{Statement.} A workshop makes chairs and tables. A chair brings a profit of $5$, a table $4$. A chair takes $1$ hour of labour and $4$ units of wood, a table $3$ hours and $1$ unit. There are $10$ hours and $13$ units of wood. Write the integer model, solve its LP relaxation, round the LP solution and find the integer optimum. Give $z_{LP}$, $z_{IP}$ and the number of chairs in the integer optimum.

\textbf{Solution.}
\begingroup\sloppy
\begin{enumerate}[label=\textbf{\arabic*.}, leftmargin=*]
  \item Let $x$ be the number of chairs and $y$ the number of tables. Maximize $5x+4y$ subject to $x+3y\le10$ (labour), $4x+y\le13$ (wood), $x,y\ge0$ and $x,y\in\mathbb Z$.
  \item Drop integrality. At the vertices $(0,0)$, $(3.25,0)$, $\bar x$ and $(0,3.3333)$ the objective is $(0,\ 16.25,\ 23,\ 13.3333)$; the best is the corner where both resources run out, $\bar x=(2.6364,\ 2.4545)$, with $z_{LP}=23$.
  \item Rounding down gives $(2,\ 2)$ with profit $18$. The nearest lattice point $(3,\ 2)$ needs $14$ units of wood, more than there is.
  \item Checking the $13$ feasible lattice points gives the integer optimum $(3,\ 1)$ with $z_{IP}=19$: three chairs and one table, not the rounded plan. The relaxation promised $23$, so the gap is $4$, and no integer plan can beat the bound.
\end{enumerate}
\endgroup
\begin{answerbox}
$z_{LP}=23$ at $\bar x=(2.6364,\ 2.4545)$; $z_{IP}=19$ at $(3,\ 1)$.
\end{answerbox}

\begin{itemize}
  \item \textbf{On/off with a fixed cost.} $x\le Mz$, $z\in\{0,1\}$, cost $Fz$: any output switches the cost on.
  \item \textbf{If A then B.} $z_A\le z_B$ for binary indicators.
  \item \textbf{At most k of n options.} $\sum_iz_i\le k$.
  \item \textbf{Either one constraint or the other.} $a^\top x\le b+M(1-z)$ and $c^\top x\le d+Mz$.
  \item \textbf{Product of binaries.} $w=z_1z_2$ becomes $w\le z_1$, $w\le z_2$, $w\ge z_1+z_2-1$.
\end{itemize}

\begin{trapbox}
Rounding the relaxed optimum is not a solution method. The rounded point may break a constraint, and even when it is feasible the true integer optimum can be several steps away, as in C1.
\end{trapbox}

\begin{notebox}{How big should M be}
Any $M$ at least as large as the largest value the variable can take is valid. Larger values keep the same integer solutions but weaken the relaxation, slow down branch and bound and, when $M$ is huge, let a solver accept $z$ a tolerance away from zero as switched off.
\end{notebox}

\problemhead{C2}{Switch a machine on or not}{On your own}
\textbf{Statement.} An order of $D=7$ units can be made on two machines. Machine $i$ costs $F_i$ to switch on and $c_i$ per unit, and makes at most $u_i$ units: $F=(20,\ 6)$, $c=(2,\ 5)$, $u=(5,\ 8)$. Write the model with binary $z_i$ and $x_i\le M_iz_i$, and find the cheapest plan. Then relax $z_i\in\{0,1\}$ to $0\le z_i\le1$ and solve the relaxation for $M_i=u_i$ and for $M_i=10\,u_i$. Give the three values.

\textbf{Solution.}
\begingroup\sloppy
\begin{enumerate}[label=\textbf{\arabic*.}, leftmargin=*]
  \item Decisions: output $x_i\ge0$ and a switch $z_i\in\{0,1\}$ for each machine. Objective $\sum_i(F_iz_i+c_ix_i)$; demand $x_1+x_2=D$; capacity $x_i\le u_i$; link $x_i\le M_iz_i$, so that output forces the fixed cost. The smallest valid choice is $M_i=u_i$, and then the link already contains the capacity.
  \item There are only four on/off patterns. Nothing on, or machine 1 alone with capacity $5< D$, cannot cover the demand. Machine 2 alone costs $F_2+c_2D=41$. Both on, with the cheap machine 1 at full capacity and the rest on machine 2, cost $46$. The integer optimum is $41$: one machine, the dearer one per unit.
  \item Relax $z_i\in[0,1]$. Each unit on machine $i$ now pays $c_i+F_i/M_i$. With $M=u$ these rates are $(6,\ 5.75)$, so all of the demand goes to machine 2 and the bound is $40.25$, just below the integer optimum.
  \item Take instead $M_i$ equal to $10$ capacities. The rates become $(2.4,\ 5.075)$: the fixed costs have almost vanished, machine 1 is filled to capacity and the bound drops to $22.15$. The integer model is the same, but the relaxation now says little about it, and a branch-and-bound solver would have to work much harder.
\end{enumerate}
\endgroup
\begin{answerbox}
Integer optimum $41$ (machine 2 alone); relaxation $40.25$ with $M=u$ and $22.15$ with the inflated $M$.
\end{answerbox}

\problemhead{C3}{A knapsack and its relaxation}{On your own}
\textbf{Statement.} Four items have values $(8,\ 6,\ 5,\ 4)$ and weights $(3,\ 5,\ 2,\ 8)$; the knapsack holds $9$. Write the model with binary $x_i$. Solve the LP relaxation by sorting the items by value per unit of weight, then find the best packing. Give the LP bound, the integer optimum and the number of items packed.

\textbf{Solution.}
\begingroup\sloppy
\begin{enumerate}[label=\textbf{\arabic*.}, leftmargin=*]
  \item Model: $z_i\in\{0,1\}$, maximize $\sum_iv_iz_i$ subject to $\sum_iw_iz_i\le W$ with $v=(8,\ 6,\ 5,\ 4)$, $w=(3,\ 5,\ 2,\ 8)$, $W=9$. By the ratio $v_i/w_i$ the order is $(1,\ 3,\ 2,\ 4)$.
  \item The relaxation takes whole items in this order while they fit and a share $0.8$ of the next one: the bound is $17.8$.
  \item Checking the subsets that fit, the best is the set $(1,\ 2)$ worth $14$; the gap to the bound is $3.8$. The greedy prefix is not necessarily the integer optimum.
\end{enumerate}
\endgroup
\begin{answerbox}
Relaxation $17.8$, integer optimum $14$ with the items $(1,\ 2)$.
\end{answerbox}

\begin{selfcheckbox}
A binary variable $z$ appears only in $x\le Mz$ and in the cost $Fz$ with $F>0$. What value does the relaxation give it?
\selfcheckanswer The smallest allowed one, $z=x/M$: the relaxation pays only the share $Fx/M$ of the fixed cost.
\end{selfcheckbox}

\end{document}
